Measure the water level $y$ from the middle between low and high tide. With high tide at $t = 0$:
\begin{align}
 y(t) &= A\cos(\omega t), \qquad A = \AF = \A, \qquad \omega = \omF = \om
\end{align}

\begin{abcliste}
\abc
The boat needs at least $h_1 = \hbO$ above low tide, that is $y \ge y_1$ with
\begin{align}
 y_1 &= \ybF = \hb-\A = \yb
\end{align}
measured from the middle. Solving $A\cos(\omega t) = y_1$ for $t$ (in radians):
\begin{align}
 t &= \tF \\
   &= \frac{1}{\om}\arccos\frac{\yb}{\A} \\
   &= \t \approx \result{\tHP{0}{2}}
\end{align}

\abc
Halfway between low and high tide the water rises or falls fastest; near low and high tide it hardly changes:
\begin{align}
 \hat v &= \vmaxF \\
   &= \A\times\om \\
   &= \vmax \approx \result{\vmhP{0}{2}}
\end{align}
\end{abcliste}
